% !TEX program = xelatex

\documentclass{resume}
%\usepackage{zh_CN-Adobefonts_external} % Simplified Chinese Support using external fonts (./fonts/zh_CN-Adobe/)
%\usepackage{zh_CN-Adobefonts_internal} % Simplified Chinese Support using system fonts
\usepackage{academicons}
\usepackage{xspace}

\newcommand{\etal}{et al.\@\xspace}

\begin{document}
\pagenumbering{gobble} % suppress displaying page number

\name{Jinguang Tong}

\basicInfo{
  \email{hirotong0715@gmail.com} \textperiodcentered\ 
  \phone{(+86) 18143465557 } % \textperiodcentered\ 
  % \linkedin[billryan8]{https://www.linkedin.com/in/billryan8}}
  \googlescholar[Google Scholar]{https://scholar.google.com/citations?user=KyaPzTEAAAAJ&hl=en}
}

\section{\faGraduationCap\ Education}
\datedsubsection{\textbf{Australian National University}, Australia}{2021 -- Present}
\textit{Ph.D in Computer Science}, expected January 2026 \\
\textbf{Advisors:} Prof.Hongdong Li, Dr.Chuong Nguyen \\
\textbf{Research Focus:} 3D Computer Vision, 3D Reconstruction, Neural Rendering
\datedsubsection{\textbf{Zhejiang University}, China}{2015 -- 2018}
\textit{Master in Optical Engineering}
\datedsubsection{\textbf{Tianjin University}, China}{2011 -- 2015}
\textit{B.S. in Electronics Engineering}

\section{\faUsers\ Experience}
\datedsubsection{\textbf{HJIMI Inc.} Beijing, China   \hspace{5em}\textbf{Machine Learning Engineer}}{2018 -- 2021}
\begin{itemize}
    \item Developed robust face detection, anti-spoofing, and recognition algorithms for a commercial face-id system.
    \item Passed the National Bank Card Testing Center’s (BCTC) security certification for the facial recognition payment system.
\end{itemize}

\section{\faFileText~Preprints}
\begin{enumerate}
    \item \textbf{Jinguang Tong}, Xuesong Li, Sundaram Muthu, Fahira Afzal Maken, Chuong Nguyen, Hongdong Li. ``S$^2$GS: Self-supervised Gaussian Segmentation for Automatic 3D Object Scanning'', Under review.
    \item Qing Zhang, Zehao Chen, \textbf{Jinguang Tong}, Jing Zhang, Jie Hong and Xuesong Li. ``Viewpoint Consistency in 3D Generation via Attention and CLIP Guidance", arXiv, 2024.
    \item Xuesong Li. \textbf{Jinguang Tong}, Jie Hong, Vivien Rolland and Lars Petersson, ``DGNS: Deformable Gaussian Splatting and Dynamic Neural Surface for Monocular Dynamic 3D Reconstruction", arXiv, 2024.
\end{enumerate}
    

\section{\faBookmark~Publications}
\begin{enumerate}
    \item \textbf{Jinguang Tong}, Xuesong Li, Fahira Afzal Maken, Sundaram Muthu, Chuong Nguyen, and Hongdong Li. ``GS-2DGS: Geometrically supervised 2DGS for reflective object reconstruction", In CVPR, 2025.
    \item Chushan Zhang*, \textbf{Jinguang Tong*}, Tao Jun Lin, Chuong Nguyen, and Hongdong Li. ``PMVC: Promoting Multi-View Consistency for 3D Scene Reconstruction", In WACV, 2024. (* equal contribution) 
    \item \textbf{Jinguang Tong}, Sundaram Muthu, Fahira Afzal Maken, Chuong Nguyen, and Hongdong Li. ``Seeing through the glass: Neural 3d reconstruction of object inside a transparent container", In CVPR, 2023.
    \item Jiahao Ma, \textbf{Jinguang Tong}, Shan Wang, Wei Zhao, Zicheng Duan, and Chuong Nguyen. ``Voxelized 3d feature aggregation for multiview detection", In DICTA, 2024.
\end{enumerate}

%%%%%%%%%%%%%%%%%%%%%%%%%%%%%%%%%%%%%%%%%%%%%%%%%%%%%%%%%%%%%%%%%%%%%%
% SKILLS & INTERESTS
%%%%%%%%%%%%%%%%%%%%%%%%%%%%%%%%%%%%%%%%%%%%%%%%%%%%%%%%%%%%%%%%%%%%%%
\section{\faCogs\ Skills \& Interests}
\subsection*{Technical Skills:}
\begin{itemize}[parsep=0.5ex]
  \item \textbf{Programming:} Python, C/C++, PyTorch
  \item \textbf{Methodologies:} Neural Rendering, 3D Gaussian Splatting, Diffusion models
\end{itemize}

\subsection*{Research Interests:}
I am deeply interested in advancing the field of 3D computer vision with a focus on both reconstruction and generation. My research explores innovative approaches in:
\begin{itemize}[noitemsep]
  \item \textbf{3D Reconstruction:} Investigating robust multi-view reconstruction methods and novel techniques for recovering accurate 3D geometry from sparse or dynamic data.
  \item \textbf{Neural Rendering:} Pushing the boundaries of neural rendering techniques to produce photorealistic images.
  \item \textbf{3D Generation:} Developing generative models for creating high-fidelity 3D contents.
\end{itemize}


% Reference Test
%\datedsubsection{\textbf{Paper Title\cite{zaharia2012resilient}}}{May. 2015}
%An xxx optimized for xxx\cite{verma2015large}
%\begin{itemize}
%  \item main contribution
%\end{itemize}

% \section{\faCogs\ Skills}
% \begin{itemize}[parsep=0.5ex]
%   \item Knowledge: NeRF, 3DGS, Neural Rendering, Diffusion model
%   \item Programming Languages: Python, C / C++
%   \item Language: English - Fluent, Mandarian - Native speaker
% \end{itemize}


\section{\faInfo\ Teaching and Service}
\subsection*{Teaching Assistant:}
\begin{itemize}[parsep=0.5ex]
  \item (Head) COMP4610, Computer Graphics, ANU, (2023S2, 2024S1).
  \item ENGN8501, Advanced Topics in Computer Vision, ANU, (2022S2, 2023S2).
  \item ENGN6528, Computer Vision, ANU, (2023S1, 2024S1).
\end{itemize}
\subsection{Professional Service:}
\begin{itemize}[parsep=0.5ex]
    \item Conference Reviewer: CVPR, ICLR.
    \item Jornal Reviewer: IEEE TPAMI, IEEE RA-L, IEEE TCSVT, IET CV.
\end{itemize}

\section{\faHourglass\ Honors and Awards}
\datedline{Outstanding Graduate Student, Zhejiang University}{March. 2018}
\datedline{Outstanding Undergraduate Student}{Jun. 2015}

%% Reference
%\newpage
%\bibliographystyle{IEEETran}
%\bibliography{mycite}
\end{document}
